% !TeX root = multianalysis.tex
% !TeX encoding = UTF-8
% !TeX program = XeLaTeX


\begin{frame}[fragile]
    \frametitle{参考文献}
    % \begin{multicols}{1}
    \tiny
    \begin{thebibliography}{99}
        \bibitem{}
        \textsc{何晓群}.
        \newblock 多元统计分析 [M]. 2019.
        \newblock 北京: 中国人民大学出版社

        \bibitem{}
        \textsc{白志东}.
        \newblock 大维统计分析 [M]. 2012.
        \newblock 北京：高等教育出版社

        % \bibitem{}
        % \textsc{陈希孺}.
        % \newblock 数理统计学教程 [M]. 2009.
        % \newblock 合肥：中国科学技术大学出版社

        % \bibitem{}
        % \textsc{陈希孺}.
        % \newblock 数理统计学简史 [M]. 2002.
        % \newblock 长沙：湖南教育出版社

        \bibitem{}
        \textsc{Vincent Spruyt}.
        \newblock A geometric interpretation of the covariance matrix [EB/OL]. 2014.
        \newblock Computer vision for dummies
        \newblock \url{https://www.visiondummy.com/2014/04/geometric-interpretation-covariance-matrix}
        
        \end{thebibliography}
    % \end{multicols}
\end{frame}


\begin{frame}
    \begin{tikzpicture}[overlay]
    \pgfmathdeclarerandomlist{symbols}{%
      {\texttt{\string Cluster}}%
      {\texttt{\string LDA}}
      {\texttt{\string PCA}}
      {\texttt{\string Factor}}}
    \foreach \y in {0,-1,...,-7} {
      \foreach \x in {0,1,...,12} {
        \pgfmathrandomitem{\randsymbol}{symbols}
        \node[black!15,rotate=45] at (\x,\y) {\randsymbol};
      }
    }
    \end{tikzpicture}
    \titlepage
\end{frame}